% GENERATED FILE. Edit canonical lessons, then run npm run generate:books.
% canonical-source-hash: fnv1a64:a4add17725eb6976
% canonical-lessons: SA-C354-nagarikah, SA-C354-janata, SA-C354-samajah, SA-C354-svatantrah, SA-C354-sanghatayati

\chapter{Citizens and Society}
\label{ch:sa-citizens-and-society}
\hlchaptermodality{354}

\begin{chapteropening}
\textbf{By the end of this chapter:} \emph{I can talk about citizens, the public and society.}

Complete the last lesson of chapter 354: I can talk about citizens, the public and society.
\end{chapteropening}

\section[\texorpdfstring{\sk{नागरिकः} (nāgarikaḥ)}{नागरिकः (nāgarikaḥ)}]{\sk{नागरिकः} (nāgarikaḥ) — a citizen (of a country)}

\label{lesson:SA-C354-nagarikah}



\begin{quote}
Before the new one: say the Sanskrit for a governor, then the Sanskrit for an official order.
\end{quote}

\subsection*{You'll want to know: \sk{नागरिकः}}
\textbf{\sk{नागरिकः}} — \emph{nāgarikaḥ} — \textquotedblleft{}a citizen (of a country)\textquotedblright{}.

\begin{quote}
\sk{नागरिकः} \sk{नियमान्} \sk{पालयेत्।}
\end{quote}

\begin{quote}
\emph{nāgarikaḥ niyamān pālayet.}
\end{quote}

\begin{quote}
A citizen should keep the laws.
\end{quote}

A masculine noun.

\begin{cousinweb}
\textbf{\sk{नागरिकः}} (\emph{nāgarikaḥ}) is built on \emph{nagara}, \textquotedblleft{}a town\textquotedblright{}; MW gives \emph{nāgarika} \textquotedblleft{}born or living in a town, civic\textquotedblright{}, and \textquotedblleft{}a citizen\textquotedblright{} is how modern Sanskrit uses it (MW).
\end{cousinweb}

\begin{grammarlens}[title={What you've built}]
One more word for giving reasons and opinions.
\end{grammarlens}

\subsection*{Your turn}
\begin{itemize}
\raggedright
  \item \emph{Say it:} \emph{nāgarikaḥ}
  \item \emph{Say it:} \emph{nāgarikaḥ}, once more
  \item \emph{Say it:} \emph{nāgarikaḥ niyamān pālayet.} — A citizen should keep the laws.
  \item \emph{Recall:} say the Sanskrit for a governor, then the Sanskrit for an official order, then say \emph{nāgarikaḥ} again
\end{itemize}

\begin{tcolorbox}[breakable,colback=teal!4,colframe=teal!35!black,title={Before you move on}]
What does \sk{नागरिकः} mean? (\textquotedblleft{}A citizen (of a country)\textquotedblright{}.) Say it once more.
\end{tcolorbox}
\section[\texorpdfstring{\sk{जनता} (janatā)}{जनता (janatā)}]{\sk{जनता} (janatā) — the public}

\label{lesson:SA-C354-janata}



\begin{quote}
Before the new one: say the Sanskrit for an official order, then the Sanskrit for a citizen (of a country).
\end{quote}

\subsection*{You'll want to know: \sk{जनता}}
\textbf{\sk{जनता}} — \emph{janatā} — \textquotedblleft{}the public\textquotedblright{}.

\begin{quote}
\sk{जनता} \sk{नूतनां} \sk{नीतिं} \sk{न} \sk{इच्छति।}
\end{quote}

\begin{quote}
\emph{janatā nūtanāṃ nītiṃ na icchati.}
\end{quote}

\begin{quote}
The public does not want the new policy.
\end{quote}

A feminine noun.

\begin{cousinweb}
\textbf{\sk{जनता}} (\emph{janatā}) is built on \emph{jana}, \textquotedblleft{}people\textquotedblright{}, with the collective ending \emph{-tā}; MW gives \emph{jana-tā} \textquotedblleft{}a number of men, an assemblage of people, a community, subjects\textquotedblright{} (MW).
\end{cousinweb}

\begin{grammarlens}[title={What you've built}]
One more word for giving reasons and opinions.
\end{grammarlens}

\subsection*{Your turn}
\begin{itemize}
\raggedright
  \item \emph{Say it:} \emph{janatā}
  \item \emph{Say it:} \emph{janatā}, once more
  \item \emph{Say it:} \emph{janatā nūtanāṃ nītiṃ na icchati.} — The public does not want the new policy.
  \item \emph{Recall:} say the Sanskrit for an official order, then the Sanskrit for a citizen (of a country), then say \emph{janatā} again
\end{itemize}

\begin{tcolorbox}[breakable,colback=teal!4,colframe=teal!35!black,title={Before you move on}]
What does \sk{जनता} mean? (\textquotedblleft{}The public\textquotedblright{}.) Say it once more.
\end{tcolorbox}
\section[\texorpdfstring{\sk{समाजः} (samājaḥ)}{समाजः (samājaḥ)}]{\sk{समाजः} (samājaḥ) — society}

\label{lesson:SA-C354-samajah}



\begin{quote}
Before the new one: say the Sanskrit for a citizen (of a country), then the Sanskrit for the public.
\end{quote}

\subsection*{You'll want to know: \sk{समाजः}}
\textbf{\sk{समाजः}} — \emph{samājaḥ} — \textquotedblleft{}society\textquotedblright{}.

\begin{quote}
\sk{समाजे} \sk{सर्वे} \sk{जनाः} \sk{समाः} \sk{सन्ति।}
\end{quote}

\begin{quote}
\emph{samāje sarve janāḥ samāḥ santi.}
\end{quote}

\begin{quote}
In society all people are equal.
\end{quote}

A masculine noun.

\begin{cousinweb}
\textbf{\sk{समाजः}} (\emph{samājaḥ}) is built on \emph{sam}, \textquotedblleft{}together\textquotedblright{}, and the root √aj, \textquotedblleft{}to drive\textquotedblright{}; MW gives \emph{sam-āja} \textquotedblleft{}a meeting, an assembly, a society, a company\textquotedblright{} (MW).
\end{cousinweb}

\begin{grammarlens}[title={What you've built}]
One more word for giving reasons and opinions.
\end{grammarlens}

\subsection*{Your turn}
\begin{itemize}
\raggedright
  \item \emph{Say it:} \emph{samājaḥ}
  \item \emph{Say it:} \emph{samājaḥ}, once more
  \item \emph{Say it:} \emph{samāje sarve janāḥ samāḥ santi.} — In society all people are equal.
  \item \emph{Recall:} say the Sanskrit for a citizen (of a country), then the Sanskrit for the public, then say \emph{samājaḥ} again
\end{itemize}

\begin{tcolorbox}[breakable,colback=teal!4,colframe=teal!35!black,title={Before you move on}]
What does \sk{समाजः} mean? (\textquotedblleft{}Society\textquotedblright{}.) Say it once more.
\end{tcolorbox}
\section[\texorpdfstring{\sk{स्वतन्त्रः} (svatantraḥ)}{स्वतन्त्रः (svatantraḥ)}]{\sk{स्वतन्त्रः} (svatantraḥ) — independent}

\label{lesson:SA-C354-svatantrah}



\begin{quote}
Before the new one: say the Sanskrit for the public, then the Sanskrit for society.
\end{quote}

\subsection*{You'll want to know: \sk{स्वतन्त्रः}}
\textbf{\sk{स्वतन्त्रः}} — \emph{svatantraḥ} — \textquotedblleft{}independent\textquotedblright{}.

\begin{quote}
\sk{भारतं} \sk{स्वतन्त्रं} \sk{राष्ट्रम्} \sk{अस्ति।}
\end{quote}

\begin{quote}
\emph{bhārataṃ svatantraṃ rāṣṭram asti.}
\end{quote}

\begin{quote}
India is an independent nation.
\end{quote}

\begin{cousinweb}
\textbf{\sk{स्वतन्त्रः}} (\emph{svatantraḥ}) is built from \emph{sva}, \textquotedblleft{}own\textquotedblright{}, and \emph{tantra}, \textquotedblleft{}a loom, a system\textquotedblright{}; MW gives \emph{sva-tantra} \textquotedblleft{}self-dependence, independence, freedom\textquotedblright{} (MW).
\end{cousinweb}

\begin{grammarlens}[title={What you've built}]
One more word for giving reasons and opinions.
\end{grammarlens}

\subsection*{Your turn}
\begin{itemize}
\raggedright
  \item \emph{Say it:} \emph{svatantraḥ}
  \item \emph{Say it:} \emph{svatantraḥ}, once more
  \item \emph{Say it:} \emph{bhārataṃ svatantraṃ rāṣṭram asti.} — India is an independent nation.
  \item \emph{Recall:} say the Sanskrit for the public, then the Sanskrit for society, then say \emph{svatantraḥ} again
\end{itemize}

\begin{tcolorbox}[breakable,colback=teal!4,colframe=teal!35!black,title={Before you move on}]
What does \sk{स्वतन्त्रः} mean? (\textquotedblleft{}Independent\textquotedblright{}.) Say it once more.
\end{tcolorbox}
\section[\texorpdfstring{\sk{संघटयति} (saṅghaṭayati)}{संघटयति (saṅghaṭayati)}]{\sk{संघटयति} (saṅghaṭayati) — organizes; brings together}

\label{lesson:SA-C354-sanghatayati}



\begin{quote}
Before the new one: say the Sanskrit for society, then the Sanskrit for independent.
\end{quote}

\subsection*{You'll want to know: \sk{संघटयति}}
\textbf{\sk{संघटयति}} — \emph{saṅghaṭayati} — \textquotedblleft{}organizes; brings together\textquotedblright{}.

\begin{quote}
\sk{मीरा} \sk{ग्रामस्य} \sk{जनान्} \sk{संघटयति।}
\end{quote}

\begin{quote}
\emph{mīrā grāmasya janān saṅghaṭayati.}
\end{quote}

\begin{quote}
Mira brings the village people together.
\end{quote}

Past (she …): \sk{सा} \sk{सभां} \sk{संघटितवती}, \emph{sā sabhāṃ saṅghaṭitavatī}.

\begin{cousinweb}
\textbf{\sk{संघटयति}} (\emph{saṅghaṭayati}) is built on \emph{sam}, \textquotedblleft{}together\textquotedblright{}, and the root √ghaṭ, \textquotedblleft{}to join\textquotedblright{}, in its causative form; MW gives \emph{saṃ-√ghaṭ} \textquotedblleft{}to cause to assemble, to collect, to join together\textquotedblright{}, and \textquotedblleft{}organizes\textquotedblright{} is its modern sense (MW).
\end{cousinweb}

\begin{grammarlens}[title={What you've built}]
One more word for giving reasons and opinions.
\end{grammarlens}

\subsection*{Your turn}
\begin{itemize}
\raggedright
  \item \emph{Say it:} \emph{saṅghaṭayati}
  \item \emph{Say it:} \emph{saṅghaṭayati}, once more
  \item \emph{Say it:} \emph{mīrā grāmasya janān saṅghaṭayati.} — Mira brings the village people together.
  \item \emph{Recall:} say the Sanskrit for society, then the Sanskrit for independent, then say \emph{saṅghaṭayati} again
\end{itemize}

\begin{tcolorbox}[breakable,colback=teal!4,colframe=teal!35!black,title={Before you move on}]
What does \sk{संघटयति} mean? (\textquotedblleft{}Organizes; brings together\textquotedblright{}.) Say it once more.
\end{tcolorbox}
